\section{Fijación de precios de los datos: los datos no se venden por GB}

A los datos se les llama el petróleo de la nueva era, pero esta metáfora puede inducir a error. El petróleo se vende por cantidad física, mientras que el valor de los datos depende de su autenticidad, escasez, verificabilidad, reutilización y beneficio marginal. Un mismo 1GB de datos puede ser una colección de páginas web repetidas o un conjunto de casos de falla clave de la cola larga; su valor es completamente distinto.

\begin{figure}[H]
\centering
\includegraphics[width=0.92\textwidth]{data-pricing.png}
\caption{Las cinco dimensiones de la fijación de precios de los datos: autenticidad, escasez, verificabilidad, reutilización y beneficio marginal.}
\end{figure}

\begin{knowledgebox}{Lectura de la figura: el precio de los datos proviene de su contribución marginal}
Los cinco factores de la figura determinan en conjunto el valor de los datos. Cuanto más auténticos, escasos y verificables sean, y cuanto mejor se transfieran, más valen; sin embargo, cuantos más datos del mismo tipo haya, menor es su beneficio marginal. Lo esencial de la industria de datos no es acumular volumen, sino encontrar datos de alto beneficio marginal.
\end{knowledgebox}

\subsection{Por qué los datos de fallas son caros}

Los datos de fallas son caros porque le indican al modelo dónde están sus límites. Muchas trayectorias exitosas pueden parecerse entre sí, mientras que las trayectorias fallidas suelen exponer defectos del sistema, escenarios de la cola larga y estrategias de recuperación. Para la conducción autónoma, los robots y los Agent, los pares de falla y corrección tienen más valor de entrenamiento que las simples demostraciones exitosas.

\begin{importantbox}{Los datos más efectivos: primero la falla, después el éxito}
Los datos que primero fallan y después tienen éxito contienen el error, el diagnóstico, la corrección y el resultado. Se acercan más al proceso de aprendizaje que los datos de puro éxito y también son más adecuados para entrenar la capacidad de recuperación.
\end{importantbox}

\subsection{Resumen de la sección}

El precio de los datos debe fijarse según su contribución marginal a la mejora de capacidades, no según su volumen. Los datos de alto valor suelen ser auténticos, escasos y verificables, y son capaces de exponer los límites del modelo.

